\chapter{编译、调试与常见错误}
\label{app:toolchain}

本附录介绍如何安装编译器、编译运行 C 程序、读懂报错，以及用调试器定位问题。不同系统细节不同，但核心命令是相通的。

\section{选择工具链}

\begin{center}
\begin{tabular}{@{}lll@{}}
\toprule
系统 & 推荐编译器 & 说明 \\
\midrule
Windows & MinGW-w64（GCC） & 免费，命令与 Linux 接近 \\
Linux & GCC / Clang & 通常已预装 \\
macOS & Clang & 安装 Xcode Command Line Tools \\
\bottomrule
\end{tabular}
\end{center}

本书推荐 \textbf{MinGW-w64 + VS Code}：编辑器轻量，GCC 的报错信息清晰，和后面的嵌入式开发环境也接近。

\section{Windows：安装 MinGW-w64}

最简单的方式是通过 MSYS2：

\begin{enumerate}
  \item 从 \url{https://www.msys2.org/} 下载并安装 MSYS2；
  \item 打开 MSYS2 UCRT64 终端，执行：
\begin{outbox}{MSYS2 命令}
pacman -S mingw-w64-ucrt-x86_64-gcc
\end{outbox}
  \item 把 \texttt{C:\textbackslash msys64\textbackslash ucrt64\textbackslash bin} 加入系统环境变量 \code{PATH}；
  \item 打开新的命令行，执行 \code{gcc --version}，能显示版本号就成功了。
\end{enumerate}

\section{第一个程序的编译与运行}

新建 \file{hello.c}：

\begin{codebox}{hello.c}
#include <stdio.h>

int main(void)
{
    printf("Hello, world!\n");
    return 0;
}
\end{codebox}

在文件所在目录打开终端，执行：

\begin{outbox}{编译与运行}
gcc hello.c -o hello -std=c89 -Wall -Wextra -pedantic
hello.exe        （Windows）
./hello          （Linux / macOS）
\end{outbox}

\begin{center}
\begin{tabular}{@{}ll@{}}
\toprule
参数 & 作用 \\
\midrule
\code{-o hello} & 指定输出文件名 \\
\code{-std=c89} & 按 C89 标准编译（本书要求） \\
\code{-Wall} & 打开常用警告 \\
\code{-Wextra} & 打开更多警告 \\
\code{-pedantic} & 严格按标准检查 \\
\code{-g} & 生成调试信息（配合 gdb） \\
\code{-lm} & 链接数学库（用了 \file{math.h} 时加在命令末尾） \\
\bottomrule
\end{tabular}
\end{center}

\begin{tipbox}[把警告当错误看]
初学阶段建议加 \code{-Wall -Wextra}。很多“能跑但结果不对”的 bug，编译器其实已经用警告提醒你了。更严格可以加 \code{-Werror}，把警告直接变成错误。
\end{tipbox}

\section{读懂错误信息}

GCC 的报错格式通常是：

\begin{outbox}{错误信息结构}
文件名:行号:列号: error: 错误描述
    源代码行
    ^
\end{outbox}

例如：

\begin{outbox}{示例}
hello.c:5:5: error: expected ';' before 'return'
     return 0;
     ^~~~~~
\end{outbox}

意思是第 5 行 \code{return} 前面少了一个分号。\textbf{先看第一条错误}，后面的错误常常是连锁反应。

\section{常见编译错误对照}

\begin{center}
\begin{tabular}{@{}p{0.42\textwidth}p{0.5\textwidth}@{}}
\toprule
错误信息 & 常见原因 \\
\midrule
\code{expected ';' before ...} & 上一行漏分号 \\
\code{expected declaration or statement at end of input} & 花括号不配对 \\
\code{'printf' undeclared} & 忘了 \texttt{\#include <stdio.h>} \\
\code{implicit declaration of function} & 函数没声明或头文件没包含 \\
\code{undefined reference to 'main'} & 没有 \code{main} 函数 \\
\code{undefined reference to 'sqrt'} & 链接时忘了 \code{-lm} \\
\code{ISO C90 forbids mixed declarations and code} & 声明写在了语句后面 \\
\texttt{stray '\textbackslash 357' in program} & 用了中文标点 \\
\code{variable-sized type used in sizeof} & C89 不允许变长数组 \\
\code{assignment makes pointer from integer} & 指针类型不匹配 \\
\bottomrule
\end{tabular}
\end{center}

\section{用 gdb 调试}

如果程序能编译但运行结果不对，或者崩溃，可以用 gdb。

\begin{outbox}{编译时加 -g，然后启动 gdb}
gcc bug.c -o bug -g -std=c89 -Wall
gdb bug.exe
\end{outbox}

常用命令：

\begin{center}
\begin{tabular}{@{}ll@{}}
\toprule
命令 & 作用 \\
\midrule
\code{run} & 运行程序 \\
\code{break main} & 在 main 函数处下断点 \\
\code{break 文件名:行号} & 在指定行下断点 \\
\code{next} & 执行下一行（不进入函数） \\
\code{step} & 执行下一行（进入函数） \\
\code{print 变量} & 打印变量的值 \\
\code{continue} & 继续运行到下一个断点 \\
\code{quit} & 退出 \\
\bottomrule
\end{tabular}
\end{center}

\begin{tipbox}[没有 gdb 怎么办]
用 \code{printf} 打印中间变量，同样能定位大部分问题。第 2 章介绍过这种方法。关键是：\textbf{在可疑的位置打印输入、中间结果和边界值}。
\end{tipbox}

\section{VS Code 配置}

安装 VS Code 和 C/C++ 扩展后，在项目目录新建 \file{.vscode/tasks.json}：

\begin{codebox}{tasks.json —— 编译任务}
{
    "version": "2.0.0",
    "tasks": [
        {
            "label": "build",
            "type": "shell",
            "command": "gcc",
            "args": [
                "-std=c89", "-Wall", "-Wextra", "-pedantic", "-g",
                "${file}", "-o", "${fileDirname}\\${fileBasenameNoExtension}.exe"
            ],
            "group": { "kind": "build", "isDefault": true }
        }
    ]
}
\end{codebox}

按 \code{Ctrl+Shift+B} 即可编译。运行可以在终端手动执行生成的 \file{.exe}。

\section{多文件项目}

\begin{outbox}{编译多个源文件}
gcc main.c math_utils.c -o app -std=c89 -Wall -Wextra
\end{outbox}

如果文件较多，可以写 Makefile 或使用 CMake。初学阶段先把“一个文件一个程序”练熟，再学分文件。

\section{编码问题}

\begin{itemize}
  \item 源文件统一保存为 \textbf{UTF-8}（无 BOM）；
  \item Windows 控制台默认可能是 GBK，输出中文可能乱码，可以执行 \code{chcp 65001} 切换到 UTF-8；
  \item 不要在代码里混用中文全角标点和英文半角标点，这是编译错误的常见来源。
\end{itemize}

\section{推荐的学习流程}

\begin{enumerate}
  \item 先手写代码，不要复制粘贴；
  \item 编译，打开所有警告；
  \item 有错先读第一条错误信息，按行号定位；
  \item 改完再编译，直到零警告；
  \item 运行，用边界值测试；
  \item 结果不对就打印中间变量或上调试器；
  \item 最后把代码整理干净，加上必要注释。
\end{enumerate}
